\chapter{Introduction}

\section{The Past, Present, and Future of Common Lisp\index{Common Lisp}}

\subsection{Lisp Yesterday}

John McCarthy\index{McCarthy, John} worked out the basic principles of Lisp in 1958, at
MIT. The name stood for ``LISt Processing'': the language was built
for manipulating symbols and lists of symbols, where the other
languages of the day were built for arithmetic. A list, in this
sense, is just a sequence of items, like a shopping list or a
checklist.

The idea that set Lisp apart was that a Lisp program is itself written
as lists. The language can therefore read, build and transform
programs with the same operations it uses on any other data. Everything
that makes Lisp unusual follows from that.

Lisp spread quickly through the research world and split into many
dialects. In the early 1980s the people behind the main ones agreed to
merge them into one language, Common Lisp. Guy Steele described it in
the book \emph{Common Lisp: the Language} in 1984, and after ten more
years of work it became an ANSI standard in 1994. It was the first
object-oriented language to be standardized by ANSI\index{ANSI standard}.

\subsection{Lisp Today}

The standard has not changed since 1994, and this turns out to be one
of the best things about the language. A Common Lisp program written
thirty years ago still runs today, usually without touching a line of
it, on implementations and machines that did not exist when it was
written. Libraries do not rot when a new version of the language
comes out, because one does not come out. Code you write this year
will not need rewriting for the language's sake.

The language can stand still because it can be extended from the
inside. Most of what other languages add in new versions (a new kind
of loop, pattern matching, a new way to manage resources) can be added
to Common Lisp by a library, and looks no different from what is built
in. Section~\ref{sec:macros} shows how.

Meanwhile the implementations and tools have kept moving. There are
several maintained implementations of Common Lisp, most of them free
and open source, and the best of them compile to efficient machine
code. Some two thousand
open-source libraries can be installed with one command. Common Lisp
is in daily commercial use in engineering design, scheduling and
logistics, finance, and research computing, and it has a small, steady
community of people who chose it on purpose.

\subsection{Lisp Tomorrow}

For most of the history of software, people have hoped for tools that
would let them describe a program at a high level and have the detailed
code written for them. Lisp's answer has always been that the
high-level description and the program should be the same thing,
written in one language. A Common Lisp program can generate other
Common Lisp code, normally at compile time, so you can build the
language up toward your problem until the program reads like a
description of it.

That matters more now, not less. A growing share of code is drafted by
AI models and then checked by people. A language with a small, regular
syntax, a precise standard, and a running image that can be asked what
is true right now is a good place to do that kind of work: the model
can try an expression and see the answer before it proposes anything,
and so can you. Section~\ref{sec:agents} says more.

Several long-standing features of Common Lisp also save time and
people:

\begin{itemize}
\item \textbf{Automatic memory management\index{garbage collection}.} You do not
allocate and free memory by hand. Lisp was the first language to have
a garbage collector.

\item \textbf{Dynamic typing\index{dynamic typing}.} Values have types, but variables
do not have to. You can write a function without declaring the type
of anything, and add declarations later where they help the compiler
or the reader.

\item \textbf{Redefinition while running.} You can add and replace
functions, classes and data in a running program without stopping it.
This is how Lisp is normally developed, and it also means that a
server can be inspected and patched while it stays up.

\item \textbf{A large standard language.} Hash tables, arbitrary-precision
integers, exact fractions, multi-dimensional arrays, a full object
system, formatted output, and a condition system for handling errors
are all part of the standard.

\item \textbf{Portability.} Because the standard is precise, code moves
between implementations and operating systems with little or no
change.
\end{itemize}

\section{Is Lisp Slow? Is It Big?}

The old complaints about Lisp were that it was too slow and needed too
much memory. They were fair on the machines of the 1970s. They have
not been true for a long time. Compiled Common Lisp runs at the speed
of other compiled languages, and a Lisp image of forty megabytes is
small next to a web browser tab.

The cost that matters today is the time of the people writing the
program. Common Lisp teams tend to be small because one person with a
good grasp of the language can build and maintain a great deal.

\section{The CL Model of Computation}

In most languages you work in a cycle: edit the source, build, run the
program from the beginning, and see what happens. In Common Lisp you
start the Lisp once and leave it running. You type an expression and
it is evaluated at once. You change a function and send just that
function to the running Lisp, which replaces the old definition and
carries on, with all its data still in memory. You test the change in
the next second.

This running Lisp is called the \emph{image}\index{image}, and the prompt where
you talk to it is called the \emph{REPL}\index{REPL}, for read-eval-print loop: it
reads an expression, evaluates it, prints the result, and waits for
the next one. Working this way changes how you write programs. You
build them in small pieces and try each piece as you go, so you find
out early when you have misunderstood something.

The next chapter shows how to set this up.
